\documentclass[10pt,a4paper]{amsart}
\usepackage[margin=19mm]{geometry}
\usepackage{amsmath,amssymb,mathtools}
\usepackage[scr=rsfs,cal=euler]{mathalfa}
\usepackage{xfrac}
\usepackage[unicode,hidelinks,bookmarks=false]{hyperref}
\newcommand{\ie}{i.e.\ }
\newcommand{\cf}{cf.\ }
\newcommand{\resp}{respectively\ }
\newcommand{\Subsection}{\subsection}
\providecommand{\nobreakdash}{\nobreak\hbox{-}}
\setlength{\emergencystretch}{2em}
\multlinegap0pt
\usepackage[shortlabels]{enumitem}
\usepackage{amssymb}
\newtheorem{theorem}{Theorem}
\title{{Global and local error bounds: characterizations via directional derivatives and tangent cones}}
\author{Yu Han}
\date{Source: arXiv:2607.29004v1 (CC BY 4.0)}
\begin{document}
\maketitle

\noindent\emph{Source and licence.} {Global and local error bounds: characterizations via directional derivatives and tangent cones}. Version arXiv:2607.29004v1 (CC BY 4.0), distributed by arXiv under the Creative Commons Attribution 4.0 International licence, http://creativecommons.org/licenses/by/4.0/, as recorded in the arXiv licence metadata for this exact version and shown on the abstract page https://arxiv.org/abs/2607.29004v1. Full licence notice: \texttt{licenses/arxiv-2607.29004-CC-BY-4.0.txt}.\par\smallskip

\noindent\textit{Editorial scope.} Counted unit: the article's theorem XLWCJK (source0.tex lines 349--387). The article's complete original proof of it is delivered with it (source0.tex lines 389--625, one \texttt{proof} environment of 237 lines). No mathematics is written, repaired or re-derived: the statement and the proof are the article's own lines, in the article's own notation, and no retained line is substituted or re-flowed.  No further source-local context is needed: every cross-reference of every retained line resolves inside the two delivered spans.  Nothing was omitted from any retained span, so the package carries no omission record.  No retained line contains a citation, so the wrapper carries no bibliography and no bibliography entry.  No retained line contains a figure environment, an image-inclusion command or drawing code, so no image is delivered and the package declares no asset dependency.  Editorial formatting only, in addition: an AMS \texttt{amsart} preamble that restates the article's own declarations and macros used by the retained lines, namely the environments \texttt{theorem} on separate counters, together with the standard packages those lines need.  The retained environments print in this document as Theorem 1 in order of appearance, whereas the article prints the same items with the numbering of its own full document.  The complete native source of the article as distributed in the arXiv e-print for this exact version is bundled unchanged as \texttt{sources/206453/source0.tex}.
\begin{theorem} \label{XLWCJK}  Consider the following statements:
	\begin{enumerate}
		\item[(i)] $\theta$ has an error bound on $A$: there exists $\tau>0$ such that
		\begin{equation}\label{EBAQJ11} 
			d(x,S_0)\le\tau \theta(x), \quad  \forall x\in A.
		\end{equation}		 
		
		\item[(ii)]  There exists   \(\kappa > 0\) such that  
		$$			d(x, S_0) \le \kappa \, \zeta(x), \quad \forall x \in X, 		$$
		where  $\zeta(x) := \inf_{a \in A} \bigl( \theta(a) + \|x - a\| \bigr)$.
		
		\item[(iii)]  There exists \(\kappa > 0\) such that  
		\begin{equation}\label{EBkkQJ} 
			d(x+e, S_0) \le \kappa (\theta(x) + \|e\|),  \quad  \forall x \in A, \;    \forall e \in X.
		\end{equation}	
		
		\item[(iv)]  There exists  $\gamma>0$ such that  
		\begin{equation}\label{EBkXJ} 
			d(x,S_0)\le \gamma\bigl(\max\{\theta(x),0\}+d(x, A)\bigr),\quad \forall x\in X .
		\end{equation}	
		
		\item[(v)] There exist \(t_0 > 0\) and \(\beta > 0\) such that 
		\[
		d_H(\Phi(t), S_0) \le \beta t,  \quad  \forall t \in [0, t_0].
		\]
		
		\item[(vi)] There exist $\kappa>0$ and an open set $U\supseteq S_0$ such that \begin{equation}\label{EBkmJ} 
			d(x,\Phi(0))\le\kappa\,d(0,\Phi^{-1}(x)),       \quad  \forall       x\in U\cap A. 	
		\end{equation}	 
	\end{enumerate}
	Then the following hold:
	\begin{enumerate}
		\item      (i) $\Leftrightarrow$ (ii), (iii) $\Rightarrow$ (i), (iv) $\Rightarrow$ (i),  (i) $\Rightarrow$ (v) and (i) $\Rightarrow$ (vi). 
		\item  If $\theta$ is Lipschitz continuous on $X$ with constant $L>0$,
		then (i) $\Rightarrow$ (iii) and (i) $\Rightarrow$ (iv).
		\item  If $A$ is bounded, then (v) $\Rightarrow$ (i).
		\item  	If $A$ is compact  and $\theta$ is lower semicontinuous on $A$, then (vi) $\Rightarrow$ (i).  
	\end{enumerate}
\end{theorem}

\begin{proof}
	\noindent\textbf{(i) $\Rightarrow$ (ii)}.
	Suppose there exists \(\tau > 0\) such that   (\ref{EBAQJ11}) is true.
	Take an arbitrary \(x \in X\). For any \(a \in A\),   it follows from  (\ref{EBAQJ11})  that
	\[
	d(x, S_0) \le \|x - a\| + d(a, S_0) \le \|x - a\| + \tau \, \theta(a).
	\]
	Set \(\kappa := \max\{\tau, 1\}\). Then  
	$  \|x - a\| + \tau \, \theta(a) \le \kappa \bigl( \|x - a\| + \theta(a) \bigr),$  and so
	$$d(x, S_0) \le \kappa (\|x - a\| + \theta(a)),  \quad  \forall a \in A.  $$
	Taking the infimum over \(a \in A\) yields  
	\[
	d(x, S_0) \le \kappa \inf_{a \in A} \bigl( \theta(a) + \|x - a\| \bigr) = \kappa \, \zeta(x).
	\]
	Since \(x \in X\) is arbitrary, this proves (ii).
	
	
	\noindent\textbf{(ii) $\Rightarrow$ (i)}.
	Suppose there exists \(\kappa > 0\) such that 
	\begin{equation}\label{ErtB1}
		d(x, S_0) \le \kappa \, \zeta(x), \quad  \forall x \in X.
	\end{equation}
	Take an arbitrary \(x \in A\). From the definition of \(\zeta\) we have  
	\[
	\zeta(x) = \inf_{a \in A} \bigl( \theta(a) + \|x - a\| \bigr) \le \theta(x) + \|x - x\| = \theta(x).
	\]
	Substituting this into    (\ref{ErtB1})   gives  
	\[
	d(x, S_0) \le \kappa \, \theta(x), \quad \forall x \in A.
	\]
	
	\noindent\textbf{(iii) $\Rightarrow$ (i)}.
	Suppose  that there exists \(\kappa > 0\) such that  (\ref{EBkkQJ}) holds.   
	Take any \(x \in A\) and set \(e = 0\). From (\ref{EBkkQJ}) we immediately obtain
	\[
	d(x, S_0) = d(x+0, S_0) \le \kappa (\theta(x) + 0) = \kappa \theta(x).
	\]
	
	
	\noindent\textbf{(iv) $\Rightarrow$ (i)}.
	Assume that  there exists $\gamma>0$ such that (\ref{EBkXJ}) is satisfied.
	Take any $x\in A$. Then $d(x, A)=0$  and   $\theta(x)\ge0$, and so $\max\{\theta(x),0\}=\theta(x)$. Substituting into (\ref{EBkXJ}) immediately yields
	\[
	d(x,S_0)\le \gamma\,\theta(x),\quad \forall x\in A.
	\]
	
	
	
	\noindent\textbf{(i) $\Rightarrow$ (v)}.
	Suppose that  there exists \(\tau > 0\) such that   (\ref{EBAQJ11})  is satisfied.
	
	Take arbitrary \(t \ge 0\) and \(x \in \Phi(t)\). Since \(x \in A\) and \(\theta(x) \le t\),   (\ref{EBAQJ11})  yields
	\[
	d(x, S_0) \le \tau \theta(x) \le \tau t.
	\]
	Because \(S_0 \subseteq \Phi(t)\), clearly \(\sup_{y \in S_0} d(y, \Phi(t)) = 0\). Hence the Hausdorff distance satisfies
	\[
	d_H(\Phi(t), S_0) = \sup_{x \in \Phi(t)} d(x, S_0) \le \tau t, \quad \forall t \ge 0.
	\]
	Taking \(t_0 = 1\) and \(\beta = \tau\), we have for all \(t \in [0, t_0]\) that \(d_H(\Phi(t), S_0) \le \beta t\). 	 		
	
	
	\noindent\textbf{(i) $\Rightarrow$ (vi)}.
	By the definition of $\Phi$, clearly $\Phi(0)=\theta^{-1}(0)\cap A= S_0$.
	For a fixed $x\in A$,
	\[
	\Phi^{-1}(x)=\{t\ge0:\theta(x)\le t\}=[\theta(x),+\infty),
	\]
	and	so $d(0,\Phi^{-1}(x))=\inf\{|t-0|:t\ge\theta(x)\}=\theta(x)$.
	Therefore,      (\ref{EBkmJ})    is equivalent to
	\begin{equation}\label{EViiiD}
		d(x, S_0)\le\kappa\,\theta(x),\quad\forall x\in U\cap A.
	\end{equation}
	If (i) holds, taking $U=X$ shows that (vi) holds.
	
	
	
	
	
	
	
	
	
	
	\noindent\textbf{(i) $\Rightarrow$ (iii)}.
	Suppose there exists \(\tau > 0\)  such that   (\ref{EBAQJ11}) holds. Take arbitrary \(x \in A\) and \(e \in X\).
	
	It follows from  \(x \in A\)  that \(d(x+e, A) \le \|x+e - x\| = \|e\|\).
	For any \(\varepsilon > 0\), choose \(a_\varepsilon \in A\) such that
	\begin{equation}\label{eTHVDD1}
		\|x+e - a_\varepsilon\| \le d(x+e, A) + \varepsilon \le \|e\| + \varepsilon,
	\end{equation}
	Since  \(\theta\) is Lipschitz continuous on \(X\) with constant \(L\),  by (\ref{eTHVDD1}), we have 
	\[
	|\theta(x+e) - \theta(a_\varepsilon)| \le L \|(x+e) - a_\varepsilon\| \le L(\|e\| + \varepsilon),
	\]
	hence
	\begin{equation}\label{eTHVDD2}
		\theta(a_\varepsilon) \le \theta(x+e) + L(\|e\| + \varepsilon).
	\end{equation}
	We conclude from  \(a_\varepsilon \in A\),  (\ref{EBAQJ11})  and  (\ref{eTHVDD2})  that
	\begin{equation}\label{eTHVDD3}
		d(a_\varepsilon, S_0) \le \tau \theta(a_\varepsilon) \le \tau\bigl(\theta(x+e) + L(\|e\| + \varepsilon)\bigr).
	\end{equation}
	It follows from  (\ref{eTHVDD1})  and  (\ref{eTHVDD3})    that
	\[
	\begin{aligned}
		d(x+e, S_0) &\le \|x+e - a_\varepsilon\| + d(a_\varepsilon, S_0) \\
		&\le (\|e\| + \varepsilon) + \tau\bigl(\theta(x+e) + L(\|e\| + \varepsilon)\bigr) \\
		&= \tau \theta(x+e) + (1 + \tau L)(\|e\| + \varepsilon).
	\end{aligned}
	\]
	Letting \(\varepsilon \to 0^+\), we obtain
	\begin{equation}\label{eq:d-xe}
		d(x+e, S_0) \le \tau \theta(x+e) + (1 + \tau L)\|e\|.
	\end{equation}
	Since  \(\theta\) is Lipschitz continuous, \(\theta(x+e) \le \theta(x) + L\|e\|\). Substituting into \eqref{eq:d-xe} yields
	\[
	\begin{aligned}
		d(x+e, S_0) &\le \tau\bigl(\theta(x) + L\|e\|\bigr) + (1 + \tau L)\|e\| \\
		&= \tau \theta(x) + (2\tau L + 1)\|e\|.
	\end{aligned}
	\]
	Due to  \(\theta(x) \ge 0\) and \(\|e\| \ge 0\), we have
	\[
	\tau \theta(x) + (2\tau L + 1)\|e\| \le (\tau + 2\tau L + 1)(\theta(x) + \|e\|).
	\]
	Setting \(\kappa = \tau + 2\tau L + 1\), we obtain
	$	d(x+e, S_0) \le \kappa (\theta(x) + \|e\|),$
	which means that  (\ref{EBkkQJ}) holds.
	
	
	
	
	\noindent\textbf{(i) $\Rightarrow$ (iv)}.
	Assume  that   there exists $\tau>0$ such that  (\ref{EBAQJ11}) holds.
	Set  
	$ \gamma := \max\{\,\tau,\;1+\tau L\,\}.$
	Let  $x\in X$      be arbitrary.   For   any    $\varepsilon>0$,  there exists $a_\varepsilon\in A$ such that
	\begin{equation}\label{EvgyB1}	\|x-a_\varepsilon\| < d(x, A)+\varepsilon  .        
	\end{equation}
	It follows from    (\ref{EBAQJ11})   and  (\ref{EvgyB1})    that   
	\begin{equation}\label{EvgyB2}
		d(x,S_0) \le \|x-a_\varepsilon\| + d(a_\varepsilon,S_0) < d(x, A)+\varepsilon + \tau\,\theta(a_\varepsilon) .
	\end{equation}
	Since  \(\theta\) is Lipschitz continuous on \(X\) with constant \(L\),  due to (\ref{EvgyB1}),  we have
	\begin{equation} \label{EvgyB3}
		\theta(a_\varepsilon) \le \theta(x) + L\|a_\varepsilon-x\| 
		< \theta(x) + L\bigl(d(x, A) + \varepsilon\bigr).  
	\end{equation}
	
	Now we distinguish two cases according to the sign of $\theta(x)$.
	
	\noindent\textbf{Case 1: $\theta(x)\ge 0$.}
	Then $\max\{\theta(x),0\}=\theta(x)$. Substituting (\ref{EvgyB3}) into (\ref{EvgyB2}),
	\[
	\begin{aligned}
		d(x,S_0) &< d(x, A)+\varepsilon + \tau\bigl(\theta(x)+L(d(x, A)+\varepsilon)\bigr) \\
		&= \tau\theta(x) + (1+\tau L)d(x, A) + (1+\tau L)\varepsilon.
	\end{aligned}
	\]
	Since $\tau\le\gamma$, $1+\tau L\le\gamma$,     $\theta(x)\ge0$  and  $  d(x, A)\ge0$, we have
	\[
	\tau\theta(x) + (1+\tau L) d(x, A) \le \gamma\theta(x) + \gamma d(x, A) = \gamma\bigl(\theta(x)+d(x, A)\bigr).
	\]
	Thus
	$	d(x,S_0) < \gamma\bigl(\theta(x)+d(x, A) \bigr) + (1+\tau L)\varepsilon.$
	
	\noindent\textbf{Case 2: $\theta(x)<0$.}
	Then $\max\{\theta(x),0\}=0$. Due to (\ref{EvgyB3})  and   $\theta(x)<0$, we have
	$\theta(a_\varepsilon) < L\bigl(d(x, A)+\varepsilon\bigr)$.
	Substituting into (\ref{EvgyB2}), we have
	\[
	\begin{aligned}
		d(x,S_0) &< d(x, A) +\varepsilon + \tau \cdot L\bigl(d(x, A) +\varepsilon\bigr) \\
		&= (1+\tau L)d(x, A) + (1+\tau L)\varepsilon.
	\end{aligned}
	\]
	Since $1+\tau L\le\gamma$ and in this case $\max\{\theta(x),0\}+d(x, A) = d(x, A)$, we obtain
	\[
	d(x,S_0) < \gamma\bigl(\max\{\theta(x),0\}+d(x, A) \bigr) + (1+\tau L)\varepsilon.
	\]
	
	Combining both cases, for any $\varepsilon>0$ we have
	\[
	d(x,S_0) < \gamma\bigl(\max\{\theta(x),0\}+d(x, A) \bigr) + (1+\tau L)\varepsilon.
	\]
	Letting $\varepsilon\to0^+$, we get
	$	d(x,S_0) \le \gamma\bigl(\max\{\theta(x),0\}+d(x, A) \bigr).	$
	This means that (\ref{EBkXJ}) is true.
	
	
	\noindent\textbf{(v) $\Rightarrow$ (i)}.
	Suppose there exist \(t_0 > 0\) and \(\beta > 0\) such that 
	\begin{equation}\label{EthhB1}
		d_H(\Phi(t), S_0) \le \beta t,  \quad  \forall t \in [0, t_0].
	\end{equation}
	
	
	Take any \(x \in A\) and set \(t = \theta(x)\). We distinguish two cases:
	
	If \(t \le t_0\), then \(x \in \Phi(t)\). By   (\ref{EthhB1}),  we have 
	\[
	d(x, S_0) \le \sup_{z \in \Phi(t)} d(z, S_0) = d_H(\Phi(t), S_0) \le \beta t = \beta \theta(x).
	\]
	
	If \(t > t_0\), since \(A\) is bounded, let \(\eta = \operatorname{diam}(A) = \sup_{a,b \in A} \|a - b\| < +\infty\). Because \(x \in A\) and \(S_0 \subseteq A\), we have \(d(x, S_0) \le \eta\). Combining this with \(t > t_0\) yields
	\[
	d(x, S_0) \le \eta < \frac{\eta}{t_0} t = \frac{\eta}{t_0} \theta(x).
	\]
	
	Setting \(\tau = \max\{\beta, \eta/t_0\}\), we obtain \(d(x, S_0) \le \tau \theta(x)\) for all \(x \in A\). 
	
	
	
	\noindent\textbf{(vi) $\Rightarrow$ (i)}.
	Assume $A$ is compact   and there exist $\kappa>0$ and an open set $U\supseteq S_0$ such that
	(\ref{EBkmJ})  is satisfied.   
	Set $H:=A\setminus U$.   Since   $A$ is compact, $U$ is open and $S_0 \subseteq U$,  we get that  $H$ is compact   and $H\cap S_0= \emptyset$.
	$\theta$ is lower semicontinuous   on the compact set $H$, hence attains its minimum $m:=\min_{x\in H}\theta(x)$.
	If $m=0$, then there exists $\bar x\in H$ with $\theta(\bar x)=0$. This yields   that   $\bar x\in S_0$, contradicting $H\cap S_0= \emptyset$. Therefore $m>0$.
	
	For any $x\in H$, since $A$ is bounded, $\beta =\operatorname{diam}(A)<\infty$, and so $d(x,S_0)\le \beta$. Together with $\theta(x)\ge m$, we have 
	\begin{equation} \label{EHYZ2}
		d(x,S_0)\le \frac{\beta}{m}\,\theta(x).  
	\end{equation}
	
	Take any $x\in A$.  If $x\in U$, noting that  (\ref{EBkmJ})  is  equivalent to (\ref{EViiiD}), we get  
	$d(x,S_0)\le\kappa\,\theta(x)$.
	
	If $x\in H$,  it follows from  (\ref{EHYZ2}) that $d(x,S_0)\le (\beta/m)\,\theta(x)$.
	Taking $\tau:=\max\{\kappa,\; \beta/m\}$, the global error bound holds:
	\[
	d(x,S_0)\le\tau\,\theta(x),\qquad \forall x\in A.
	\]
	This completes the proof. 
\end{proof}

\end{document}
